% ── §R2: Boolean Circuit Diagrams (TikZ) ──
\section{Boolean Circuit Representations}
\label{sec:circuit-diagrams}

\subsection{Level-1 Universal Contact Circuit}

The trained Level-1 circuit encodes contact propensity as a sum of prime
implicants over physicochemical group identities and separation bins.
Figure~\ref{fig:circuit-l1} depicts the essential implicants as a two-level
AND--OR network using standard rectangular gate notation.

\begin{figure}[ht]
	\centering
	\begin{tikzpicture}[
		box/.style={draw, thick, minimum width=18mm, minimum height=10mm,
				inner sep=2pt, font=\small\bfseries, fill=blue!10},
		orbox/.style={draw, very thick, minimum width=14mm, minimum height=24mm,
				inner sep=2pt, font=\small\bfseries, fill=orange!15},
		outbox/.style={draw, very thick, minimum width=20mm, minimum height=8mm,
				inner sep=2pt, font=\small\bfseries, fill=green!15},
		input/.style={font=\scriptsize, anchor=east, align=right},
		wire/.style={-{Stealth[length=5pt]}, thick},
		scale=0.9, transform shape]

		% ── AND gates as rectangles ──
		\node[box] (and1) at (0,3) {AND};
		\node[font=\tiny, anchor=south] at (and1.north) {hydrophobic $\times$ aromatic};

		\node[box] (and2) at (0,1) {AND};
		\node[font=\tiny, anchor=south] at (and2.north) {aromatic $\times$ any};

		\node[box] (and3) at (0,-1) {AND};
		\node[font=\tiny, anchor=south] at (and3.north) {sulfur $\times$ hyd/aro};

		% ── OR gate ──
		\node[orbox] (or1) at (4,1) {OR};

		% ── Output ──
		\node[outbox] (out) at (7.5,1) {CONTACT};

		% ── Input labels ──
		\node[input, anchor=east] at (-2.5,3.3) {$g_i \in \{$hyd,aro$\}$};
		\node[input, anchor=east] at (-2.5,2.8) {$g_j \in \{$hyd,aro$\}$};
		\node[input, anchor=east] at (-2.5,2.3) {$s \in [6,21)$};
			\draw[wire] (-2.4,3.0) -- (and1.west);

			\node[input, anchor=east] at (-2.5,1.3) {$g_i = $ arom};
			\node[input, anchor=east] at (-2.5,0.8) {$s \in [6,11)$};
		\draw[wire] (-2.4,1.0) -- (and2.west);

		\node[input, anchor=east] at (-2.5,-0.7) {$g_i = $ sulfur};
		\node[input, anchor=east] at (-2.5,-1.2) {$g_j \in \{$hyd,aro$\}$};
		\draw[wire] (-2.4,-1.0) -- (and3.west);

		% ── Wires from ANDs to OR ──
		\draw[wire] (and1.east) -- (or1.west);
		\draw[wire] (and2.east) -- (or1.west);
		\draw[wire] (and3.east) -- (or1.west);

		% ── Wire from OR to output ──
		\draw[wire] (or1.east) -- (out.west);

	\end{tikzpicture}
	\caption{
		\textbf{Simplified Level-1 universal contact circuit.}  Each AND block
		corresponds to one essential prime implicant discovered by Quine--McCluskey
		minimisation; the OR block combines them into the final sum-of-products
		output.  Only the three highest-scoring essential rules are shown for clarity;
		the complete circuit contains 20--22 essential implicants per fold.  Input
		labels denote physicochemical group memberships ($g$) and separation bins
		($s$).  Adapted from the trained model on PSICOV150 (holdout split).
	}
	\label{fig:circuit-l1}
\end{figure}

\subsection{Interpretation of Essential Rules}

The three implicants shown in Figure~\ref{fig:circuit-l1} correspond to
well-characterised structural interactions:
\begin{enumerate}[nosep]
	\item \textbf{AND$_1$} (hydrophobic/aromatic $\times$ hydrophobic/aromatic,
	      separation 6--21): hydrophobic-core packing between bulky side chains.
	      This rule fires on 53 on-set cells across the training set and is
	      essential in all five cross-validation folds.
	\item \textbf{AND$_2$} (aromatic $\times$ any group, separation 6--11):
	      $\pi$-stacking or edge-to-face aromatic contacts in structured regions,
	      reflecting the propensity of phenylalanine, tyrosine, and tryptophan
	      to form specific geometric arrangements.
	\item \textbf{AND$_3$} (sulfur $\times$ hydrophobic/aromatic, separation
	      6--11): cysteine-mediated thiol contacts and methionine sulfur--$\pi$
	      interactions, consistent with the known role of sulfur-containing
	      residues in stabilising local structural motifs.
\end{enumerate}

At Level-2 resolution, the circuit specifies concrete residue-pair vocabularies
such as $\{M,F\} \times \{M,F,Y,W\}$ (methionine--phenylalanine packing),
$\{I,V\} \times \{I,V,F,W\}$ (branched-chain aliphatic clustering), and
$\{I,F,C\} \times \{F,W\}$ (aromatic--sulfur contacts), all at separation
bin~0 (residues 6--10 apart).  These are precisely the residue combinations
that dominate the buried cores of globular proteins, rediscovered here by
Boolean minimisation without any biophysical prior knowledge.
